\E{icarophily}{ˌɪkəˈrɒfɪli}{n.}{Gk. Ἴκαρος Icarus + φιλία love} The love of flying too high; the disposition which reads the fable of Icarus as a complaint against wax rather than against ambition. \ex{Icarophily answers the myth with a better adhesive.} \src{Ep. Fabr.} iv. 2. --- \dv{icarophile} \textit{n.}; \dv{icarophilic} \textit{adj.} Cf. \xr{hybrionomy}.

\E{ichnomachy}{ɪkˈnɒməki}{n.}{Gk. ἴχνος track, footprint + μάχη battle} The war of the leaderboards; the contest of laboratories for the top line of a benchmark, conducted in tables and won by decimals. \sn{1} The contest itself. \sn{2} \lbl{derog.} The overfitting of a model to its examination. \ex{The ichnomachy was settled on Tuesday and reopened on Wednesday.} \src{Agora} --- \dv{ichnomachist} \textit{n.}

\E{iconoplast}{aɪˈkɒnəplæst}{n.}{Gk. εἰκών image + πλάστης moulder} \sn{1} A maker of images by machine; a model that dreams pictures from words. \sn{2} \lbl{transf.} One who shapes a movement by its images; a meme-smith. --- \dv{iconoplasty} \textit{n.} Cf. \xr{memetaxis}.

\E{ignorabimist}{ˌɪɡnɒˈræbɪmɪst}{n.}{L. \la{ignorabimus} we shall not know} One who declares that some question will never be answered, and thereby excuses the institution from trying. \ex{Every ignorabimist is refuted by a graduate student within the century.} \src{Hom. Grad.} xi. --- \dv{ignorabimism} \textit{n.}

\E{impendist}{ɪmˈpɛndɪst}{n.}{L. \la{impendere} to expend, lay out} A seed investor; one who lays out money upon a founder before there is a product, a market or, in some cases, a name. --- \dv{impendence} \textit{n.} the act of early investment. Cf. \xr{potentiarch}.

\E{incunabulant}{ˌɪnkjʊˈnæbjʊlənt}{adj.}{L. \la{incunabula} swaddling-clothes, cradle} In the cradle; said of a venture in its first months, when it occupies a garage, a spare room or a single shared document. \ex{All the great firms were once incunabulant and ridiculous.} \src{Chron. Fabr.} --- \dv{incunabulance} \textit{n.}

\E{ionotheca}{ˌaɪənəʊˈθiːkə}{n., pl. -cae}{Gk. ἰόν (neut. of ἰών going), ion + θήκη case, store} A battery; a casket in which ions are kept until they are wanted. \ex{The age of the ionotheca made the sun a night-worker.} \src{Tract. Cop.} viii. --- \dv{ionothecal} \textit{adj.}; \dv{ionothecary} \textit{n.} a maker of batteries.

\E{isegorist}{aɪˈsɛɡərɪst}{n.}{Gk. ἰσηγορία equal right of speech} A partisan of open speech upon the Agora, who would rather answer a bad post than remove it. Opp. \xr{hesychocrat}. --- \dv{isegorism} \textit{n.}

\E{jacturophobia}{ˌdʒæktjʊrəˈfəʊbiə}{n.}{L. \la{jactura} a throwing overboard, loss + Gk. φόβος fear} The incumbent's dread of loss, which outweighs any hope of gain; the reason established firms do not build what would cannibalise them. \ex{The giant died of jacturophobia, clutching its margins.} \src{Sat. Decel.} x. --- \dv{jacturophobe} \textit{n.}

\E{janitocracy}{ˌdʒænɪˈtɒkrəsi}{n.}{L. \la{janitor} doorkeeper + Gk. κράτος rule} Rule by gatekeepers; an order in which nothing new may enter without the consent of those who profit from keeping it out. --- \dv{janitocrat} \textit{n.} Cf. \xr{bradycrat}.

\E{jocularch}{ˈdʒɒkjʊlɑːk}{n.}{L. \la{jocularis} jesting + Gk. ἀρχός leader} A leader by jest; an anonymous poster whose images command a faction more surely than any manifesto. \ex{The jocularch posted one picture and the doctrine was settled.} \src{Agora} --- \dv{jocularchy} \textit{n.}

\E{kentrocracy}{kɛnˈtrɒkrəsi}{n.}{Gk. κέντρον centre + κράτος rule} Centralised control of a technology by a few laboratories or states. Opp. \xr{anacalyptism}. --- \dv{kentrocrat} \textit{n.}; \dv{kentrocratic} \textit{adj.}

\E{keraunurgy}{kəˈrɔːnɜːdʒi}{n.}{Gk. κεραυνός thunderbolt + -ουργία working} The making of thunderbolts; pulsed power, plasma confinement and the other arts by which fusion is coaxed from a machine. --- \dv{keraunurge} \textit{n.} Cf. \xr{symphysergy}.

\E{kerdophily}{kɜːˈdɒfɪli}{n.}{Gk. κέρδος gain, profit + φιλία love} The love of profit, not as greed, but as a signal that something useful has been made. \ex{Kerdophily is only the market saying thank you.} \src{Tract. Cop.} ix. --- \dv{kerdophile} \textit{n.}

\E{kleptonoesis}{ˌklɛptəʊnəʊˈiːsɪs}{n.}{Gk. κλέπτειν to steal + νόησις thought} The theft of thought; the distillation of one model's capacities into another by imitating its outputs. --- \dv{kleptonoetic} \textit{adj.}

\E{kolossurgy}{kəˈlɒsɜːdʒi}{n.}{Gk. κολοσσός colossus + -ουργία working} The building of colossi; megaprojects such as dams, canals, launch-towers and gigawatt campuses. \ex{An age without kolossurgy is an age that has forgotten its own size.} \src{Lib. Ascens.} vi. --- \dv{kolossurge} \textit{n.}

\E{kryptonoia}{ˌkrɪptəˈnɔɪə}{n.}{Gk. κρυπτός hidden + νοῦς mind} The hidden reasoning of a model, which proceeds beneath its stated answer. --- \dv{kryptonoic} \textit{adj.}

\E{kybernetrice}{ˌkaɪbɜːnəˈtriːs}{n.}{Gk. κυβερνήτης steersman + Fr. fem. -trice} The governing intelligence of a robot; the steerswoman within the machine. \ex{Give the arm a kybernetrice and it becomes a hand.} \src{Ep. Fabr.} vii. Cf. \xr{automaturgy}.

\E{kymatist}{ˈkaɪmətɪst}{n.}{Gk. κῦμα wave + -ist} One who rides waves of enthusiasm, entering each fashion at the crest and leaving before the trough. --- \dv{kymatism} \textit{n.}

\E{laboratrice}{ləˌbɒrəˈtriːs}{n.}{L. \la{laboratorium} + Fr. fem. -trice} An automated laboratory; a self-driving lab in which robots run experiments that models design. \ex{The laboratrice ran ten thousand trials while its director slept.} \src{Chron. Fabr.}

\E{lamprotype}{ˈlæmprətaɪp}{n.}{Gk. λαμπρός bright + τύπος stamp, model} A shining prototype; the demonstration that wins the investment before the product exists. \sn{1} A working demo. \sn{2} \lbl{derog.} A demo that works only once, on stage. --- \dv{lamprotypal} \textit{adj.}

\E{laterculist}{ləˈtɜːkjʊlɪst}{n.}{L. \la{laterculus} a small brick} A layer of bricks; one who prefers infrastructure to applications. \ex{The laterculist builds the road; others build the shops along it.}

\E{lexiphage}{ˈlɛksɪfeɪdʒ}{n.}{Gk. λέξις word, speech + -φάγος eating} A word-eater; a model trained upon the whole written record of mankind. \ex{The lexiphage ate the libraries and asked for more.} \src{Chron. Fabr.} --- \dv{lexiphagy} \textit{n.}

\E{limenarch}{laɪˈmiːnɑːk}{n.}{Gk. λιμήν harbour + ἀρχός ruler} The master of a spaceport; the official charged with keeping launch cadence. Cf. \xr{ouranobasis}. --- \dv{limenarchy} \textit{n.}

\E{loimophobia}{ˌlɔɪməˈfəʊbiə}{n.}{Gk. λοιμός plague + φόβος fear} The fear of engineered pestilence; the particular dread that biological tools, once open, will be turned to slaughter. Distinguished by the e/acc from mere \xr{noöphobia} and treated with more respect. --- \dv{loimophobe} \textit{n.}

\E{lucrifuge}{ˈluːkrɪfjuːdʒ}{n. \& adj.}{L. \la{lucrum} gain + \la{fugere} to flee} One who flees from profit; one who holds that a good work is tainted by being paid for. Opp. \xr{kerdophily}.

\E{machinodicy}{ˌmækɪˈnɒdɪsi}{n.}{Gk. μηχανή machine + δίκη justice, after \textit{theodicy}} The justification of the ways of machines to men; the reply to every charge laid against automation. Cf. \xr{thermodicy}. --- \dv{machinodicean} \textit{adj.}

\E{megathesaurism}{ˌmɛɡəθɪˈsɔːrɪzəm}{n.}{Gk. μέγας great + θησαυρός treasure} The raising of vast funds for a single venture; the doctrine that some goals can be reached only by the treasury of a small nation. \ex{Megathesaurism is merely arithmetic, applied boldly.} \src{Tract. Cop.} x. --- \dv{megathesaurist} \textit{n.}

\E{metagrammatist}{ˌmɛtəˈɡræmətɪst}{n.}{Gk. μεταγράφειν to rewrite + γράμμα letter} A gene-editor; one who rewrites the letters of the genome as a scribe corrects a manuscript. --- \dv{metagrammatism} \textit{n.} Cf. \xr{biotachy}.

\E{metoikist}{mɛˈtɔɪkɪst}{n.}{Gk. μέτοικος resident alien} A foreign founder or engineer settled where the building is; one admitted to a city for what they can make rather than where they were born. --- \dv{metoikism} \textit{n.}

\E{mimesiphile}{mɪˈmiːsɪfaɪl}{n.}{Gk. μίμησις imitation + φίλος loving} A fast follower; a firm that waits for another to discover a thing and then perfects it. --- \dv{mimesiphily} \textit{n.}

\E{misthophore}{ˈmɪsθəfɔː}{n.}{Gk. μισθοφόρος receiving wages} A hired builder; an engineer who is paid but does not own a share. Opp. \xr{neoctist}.

\E{mnemoclast}{ˈniːməʊklæst}{n.}{Gk. μνήμη memory + κλάστης breaker} An event that wipes a model's context or a society's memory. --- \dv{mnemoclasm} \textit{n.}

\E{mnemotheca}{ˌniːməʊˈθiːkə}{n., pl. -cae}{Gk. μνήμη memory + θήκη case, store} The memory store of a reasoning engine; its context and caches, considered as its short-term mind. \ex{Widen the mnemotheca and the model remembers the whole book.} \src{Ep. Fabr.} ix.

\E{multiplicatrix}{ˌmʌltɪplɪˈkeɪtrɪks}{n.}{L. \la{multiplicare} to multiply + -trix} Compound growth; the force by which small advantages become large ones. \ex{The multiplicatrix is patient; the decel is not.} \src{Hom. Grad.} x.

\E{nanurgy}{ˈnænɜːdʒi}{n.}{Gk. νᾶνος dwarf + -ουργία working} The working of matter at the scale of molecules; nanotechnology. --- \dv{nanurge} \textit{n.}

\E{nemesiophobe}{ˌnɛmɪsiəʊˈfəʊb}{n.}{Gk. Νέμεσις retribution + φόβος fear} One who fears that success will be punished by the gods; a doomer of the classical school. Cf. \xr{hybrionomy}.

\E{neoterist}{niˈɒtərɪst}{n.}{Gk. νεωτερισμός innovation, revolution} An innovator; one who introduces new things and is therefore suspected of disorder. --- \dv{neoterism} \textit{n.}

\E{neurosyndesis}{ˌnjʊərəʊsɪnˈdiːsɪs}{n.}{Gk. νεῦρον nerve + σύνδεσις binding together} The binding of nerve to machine; the brain-computer interface. \ex{Neurosyndesis will make the keyboard a museum piece.} \src{Lib. Ascens.} viii. --- \dv{neurosyndetic} \textit{adj.}

\E{nomothetrice}{ˌnɒməθɪˈtriːs}{n.}{Gk. νομοθέτης lawgiver + Fr. fem. -trice} A model that drafts laws. \ex{The nomothetrice wrote the statute in an hour; the committee took a year to read it.} \src{Sat. Decel.} xi.

\E{nosoclast}{ˈnɒsəʊklæst}{n.}{Gk. νόσος disease + κλάστης breaker} A breaker of diseases; a founder or scientist in biotechnology. Cf. \xr{iatrotachy}.

\E{nucleurge}{ˈnjuːkliɜːdʒ}{n.}{L. \la{nucleus} kernel + Gk. -ουργός worker} An engineer of the atomic nucleus, whether by fission or fusion. Cf. \xr{schizergy}, \xr{symphysergy}. --- \dv{nucleurgy} \textit{n.}

\E{ochetergy}{ɒˈkɛtɜːdʒi}{n.}{Gk. ὀχετός conduit, channel + ἔργον work} The carrying of energy along conduits, as by pipelines and high-voltage lines; the arteries of the \xr{diktyergy}. \ex{Generation is easy; ochetergy is where projects go to die.} \src{Tract. Cop.} xi.

\E{oistrophile}{ˈɔɪstrəfaɪl}{n.}{Gk. οἶστρος gadfly, frenzy + φίλος loving} A lover of the gadfly role; a contrarian of the Agora who stings consensus for sport. --- \dv{oistrophily} \textit{n.}

\E{omphalopolis}{ˌɒmfəˈlɒpəlɪs}{n.}{Gk. ὀμφαλός navel + πόλις city} The navel-city; the single place to which the builders of an age must come. \ex{Every omphalopolis believes itself eternal, and none has been.} \src{Chron. Fabr.}

\E{onkophily}{ɒŋˈkɒfɪli}{n.}{Gk. ὄγκος bulk, mass + φιλία love} The love of large models; the preference for scale over cleverness. Cf. \xr{scalist}. --- \dv{onkophile} \textit{n.}

\E{opiferous}{əʊˈpɪfərəs}{adj.}{L. \la{opifer} bringing aid} Bringing help; said of technologies whose benefits fall chiefly on the poor. --- \dv{opiferousness} \textit{n.}

\E{orbiculist}{ɔːˈbɪkjʊlɪst}{n.}{L. \la{orbiculus} small circle, after \textit{orbit}} One who lives or works in orbit. --- \dv{orbiculism} \textit{n.} Cf. \xr{apoikiast}.

\E{orygmatist}{ɒˈrɪɡmətɪst}{n.}{Gk. ὄρυγμα excavation, trench} A miner of asteroids. \ex{The orygmatist sees in every rock a city's worth of iron.} \src{Lib. Ascens.} ix.

\E{otiarch}{ˈəʊʃiɑːk}{n.}{L. \la{otium} leisure + Gk. ἀρχός ruler} The ruler of a post-scarcity leisure society, whose task is to keep the idle from boredom. --- \dv{otiarchy} \textit{n.} Cf. \xr{plethocracy}.
